\zlabel{triptych:concise:themes:start}
\section{The Propers: Themes and Movement}

The Son's feast calls a people into a love that must become their way of
life. God's invitation precedes their achievement, yet neither entering
the hall nor receiving his gifts makes conversion unnecessary.

\subsection{The promise calls for an answer}

The Introit begins with the Lord's saving identity and enduring lordship.
The people do not create their relationship with him by first becoming
worthy of notice. They call in trouble because he promises to hear.
The psalm verse immediately gives them something to do: listen.
Psalm 77(78)'s remembered gifts and failures make that summons searching.
Deliverance and nourishment can be received without a steadfast heart.
God's generosity makes a response possible; it does not render every
response identical.

The Gospel gives this relation a narrative shape. The feast is already
prepared when servants go to summon its guests. Refusal cannot be blamed
on the king's failure to provide. Some prefer their ordinary affairs;
others attack the messengers. After judgment on the murderers, the
invitation goes out again and gathers bad and good. The full hall displays
the king's generosity, but its mixed company still awaits his examination.
The question has moved from why someone refuses to enter to how someone
lives who has entered.

Augustine and Gregory identify the wedding garment as charity. Their
reading asks what makes presence at the feast fruitful: love appropriate
to the Bridegroom, rather than confidence in the mere fact of admission.
The Collect places the necessary change within prayer. Body and mind
are to be freed from obstruction for God's service. Freedom has this
positive end; relief from difficulty is not the whole good sought.
The people ask to become available to the Lord whose promise they have
heard.\footnote{Augustine, Sermon 90.4--6; Gregory, \work{Homilies on
the Gospels} II.38.9--12; Schuster, \work{The Sacramentary} III,
pp.~171--172.}

\subsection{The answer takes the form of a common life}

The Epistle makes the garment's moral demand visible without reducing
love to appearances. A renewed person tells the truth to those who depend
on him, refuses to cultivate resentment, and works so that a person in
need can receive help. Paul grounds this conduct in membership: one
person's speech and actions affect the body to which he belongs. Renewal
is therefore more than a private improvement whose consequences end at
the self.

The Epistle begins with renewal of the mind and ends with a gift to
need. These ends belong together. A changed way of thinking must reach
another person's material good; a visible action must arise from a
life made new. The Gospel's scrutiny presses the same depth from the
other direction, passing beyond the guests' presence to the garment.
An assembled crowd and a busy pair of hands are both capable of
concealing refusal. The Son's invitation and Paul's command call
for the person, not merely the appearance of a fitting response.

Chrysostom follows the passage from the general command to put on the
new person into these particular actions. Bellarmine's exposition of
the Communion psalm supplies their dependence on grace. Obedience is
real, but it does not earn the first gift of justification. The person
already helped by God must still ask for directed ways. This second
emphasis asks how grace changes ordinary powers of action; the garment
reading asks what gives those actions their life. Neither can dispense
with the other. Charity needs truthful and generous expression; action
needs a love directed beyond its own success.\footnote{Chrysostom,
\work{Homilies on Ephesians} XIII--XIV, on 4:23--28; Bellarmine,
\work{Commentary on the Psalms} 118:4--8.}

\newpage
\subsection{The people offer and still need help}

The Gradual draws speech and hands into prayer. Its incense and evening
sacrifice are images of an offering sought to be acceptable before God.
Augustine reads Christ's voluntary Passion in the evening sacrifice:
the Church prays in her Head, rather than establishing a separate access
to God by the force of her own achievement. Bellarmine attends to prayer's
direction. Intention can rise toward God or turn aside toward admiration.
The hands that give and the hands that pray both belong to a person whose
purpose needs purification.

The Alleluia gives speech its outward movement. The whole psalm remembers
God's covenant and saving deeds; its appointed verse commands praise,
invocation and proclamation among nations. Augustine hears the evangelists
in that announcement. Bellarmine describes love's wish that others should
know the beloved. A community called to truthful speech is also called
to generous witness. It praises a goodness whose communication brings
others joy, and even this praise requires God's assistance.

The psalm's remembered deeds culminate in the observance of God's
statutes. Praise therefore looks back to what the Lord has done and
forward to the life his gifts make possible. Its covenant, nourishment
and land are neither an empty refrain nor an escape from the demands
of the Epistle. They give thanksgiving its content. The guest who
receives and the worker who gives inhabit a history of divine
generosity. Their response has an object beyond their own performance:
the Lord whose deeds they announce.\footnote{Ps 104(105), especially
verses 8--45, supplies the context of the appointed verse 1.}

The Offertory keeps that worship within a world of actual tribulation.
God's hand preserves the singer while hostility remains. The Collect's
petition for freedom and the Offertory's confidence in trial can therefore
belong to the same worshipper. Deliverance is asked for; fidelity is
possible while the trial continues. Augustine's account of life secure
beyond an enemy's power gives charity a ground firmer than the hope of
never being hurt. The enemy can be opposed without making hatred the
source of one's safety.\footnote{Augustine, \work{Enarrationes in
Psalmos} 140.2--4; 104.1--3; 137.10--13; Bellarmine,
\work{Psalms} 140:1--2 and 104:1--2.}

\subsection{The feast heals the people it examines}

The Secret asks that the gifts offered before divine majesty be saving
for their participants. Schuster distinguishes the Eucharist's saving
power from its fruit in a recipient. The petition neither casts doubt
on Christ's gift nor makes the recipient's disposition irrelevant.
Its concern meets the Gospel at the point where the guest is already
inside. Sacred participation is a good received with responsibility,
not a certificate of completed holiness.

The Communion then puts command and petition beside one another.
God has commanded observance; the speaker wishes his own ways to be
directed accordingly. Augustine hears dependence in that wish,
Bellarmine obedient growth, Schuster delight in God's service. These
accents answer different temptations: self-sufficiency, passivity and
the belief that obedience can only be endured. A person can love the
command and still need help to live it.

The Postcommunion asks that healing reach the inclinations from which
actions proceed, freeing the people from perversity and keeping them
attached to God's commands. Its request for continuing adherence answers
the Gospel's judgment with a petition for perseverance. Gregory's
distinction between imperfect charity and charity wholly absent gives
that petition hope. The feast's medicine is needed by people whose
love is real and still growing. Their final good is the perfected
communion toward which the present gathering and its renewed life
are directed.\footnote{Schuster, \work{Sacramentary} III, pp.~173--174;
Augustine, \work{Enarrationes} 118, Aleph exposition, 4--5;
Bellarmine, \work{Psalms} 118:4--8; Gregory, \work{Gospel Homilies}
38.12--14.}
\zlabel{triptych:concise:themes:end}
\clearpage
